\section{Experiments}
\label{sec:experiments}

\subsection{Proposed model}
The proposed model is MSTNet: one STNet trunk for Sentinel-2, one for daily Xavier climate (precipitation, ETo, Rs, Tmax, and Tmin), and late fusion of MapBiomas Solo properties in the decoder.
Each trunk uses one Transformer layer with $d_{\mathrm{model}}=128$, four attention heads, $d_{\mathrm{inner}}=64$, and dropout $0.4$.
Training fits municipality-mean Sentinel-2 series; reported metrics apply the selected weights pixel by pixel and average the scores to municipalities (Section~\ref{sec:pixel-formulation}).

\subsection{Year leave-one-out design}
\label{sec:loo}
We evaluate temporal generalization with year leave-one-out (LOO) over harvest years $\{2020,2021,2022,2023,2024\}$.
For each fold with holdout year $Y$:
\begin{itemize}
  \item all municipality--years with harvest year $\neq Y$ are used for training;
  \item year $Y$ is held out entirely (no municipality from $Y$ enters training);
  \item model selection within the fold uses validation $R^2$ on year $Y$;
  \item the reported fold score is the pixel-mean municipal evaluation on year $Y$ only.
\end{itemize}
Input standardization is refit on each fold's training years.
This yields five independently trained checkpoints, one per left-out year.
Main-text tables and figures report only those LOO-year scores.
Appendix~\ref{app:muni-mean-loo} records the municipal-mean evaluation path on the same LOO years.

\subsection{Evaluation protocol}
For each LOO fold we compare municipal aggregates of pixel predictions with PAM labels on the held-out year.
We report per-fold RMSE, MAE, and $R^2$, predicted-versus-observed scatterplots for every LOO year, a municipal error map for a non-drought fold (2021), and a qualitative intramunicipal map for Guarapuava.
